\documentclass{article}
\usepackage{amsmath,amssymb,amsthm}
\usepackage{algorithm}
\usepackage{algorithmic}
\usepackage{bm}
\usepackage{enumitem}
\usepackage{hyperref}
\usepackage{geometry}
\geometry{margin=1in}
\newtheorem{theorem}{Theorem}
\newtheorem{lemma}{Lemma}
\newtheorem{assumption}{Assumption}
\newcommand{\EE}{\mathbb{E}}
\newcommand{\RR}{\mathbb{R}}
\newcommand{\norm}[1]{\left\|#1\right\|}
\newcommand{\ip}[2]{\left\langle #1,#2\right\rangle}
\newcommand{\argmin}{\operatorname*{arg\,min}}
\begin{document}

\section*{AdamW (Adam with Decoupled Weight Decay)}

\section*{Mathematical Formulation}
Let $f:\RR^{d}\rightarrow\RR$ be the (possibly non‑convex) loss function.  
At iteration $t\ge 1$ the algorithm maintains

\begin{align}
g_t &:= \nabla f(w_t; \xi_t), \qquad\text{(stochastic gradient)}\\[4pt]
m_t &= \beta_1 m_{t-1} + (1-\beta_1) g_t, \label{eq:m}\\[4pt]
v_t &= \beta_2 v_{t-1} + (1-\beta_2) g_t^{\odot 2}, \label{eq:v}
\end{align}
where $g_t^{\odot 2}$ denotes the element‑wise square and $\beta_1,\beta_2\in[0,1)$.

Bias‑corrected moments:
\begin{align}
\hat m_t &= \frac{m_t}{1-\beta_1^{t}}, \qquad
\hat v_t = \frac{v_t}{1-\beta_2^{t}}. \label{eq:bias}
\end{align}

Learning rate schedule $\{\eta_t\}_{t\ge1}$ (e.g. $\eta_t=\eta/\sqrt{t}$) and weight‑decay coefficient $\lambda\ge0$ are given.  
The parameter update of AdamW is

\begin{align}
w_{t+1}
&= w_t - \eta_t\Bigl( \frac{\hat m_t}{\sqrt{\hat v_t}+\epsilon} + \lambda w_t \Bigr), 
\qquad \epsilon>0. \label{eq:update}
\end{align}
Equation \eqref{eq:update} corresponds to applying weight decay \emph{after} the adaptive step, i.e. the decay term is \emph{decoupled} from the moment estimates.

\section*{Pseudocode}
\begin{algorithm}[H]
\caption{AdamW}
\begin{algorithmic}[1]
\REQUIRE Initial point $w_0\in\RR^{d}$, step size $\eta>0$, decay $\lambda\ge0$, 
        $\beta_1,\beta_2\in[0,1)$, $\epsilon>0$, schedule $\{\eta_t\}$.
\STATE $m_0\gets 0$, $v_0\gets 0$.
\FOR{$t=1,2,\dots,T$}
    \STATE Sample a mini‑batch $\xi_t$ and compute stochastic gradient $g_t=\nabla f(w_t;\xi_t)$.
    \STATE $m_t\gets \beta_1 m_{t-1}+(1-\beta_1)g_t$.
    \STATE $v_t\gets \beta_2 v_{t-1}+(1-\beta_2)g_t^{\odot 2}$.
    \STATE $\hat m_t\gets m_t/(1-\beta_1^{t})$, \quad $\hat v_t\gets v_t/(1-\beta_2^{t})$.
    \STATE $w_{t+1}\gets w_t-\eta_t\bigl(\hat m_t/(\sqrt{\hat v_t}+\epsilon)+\lambda w_t\bigr)$.
\ENDFOR
\RETURN $w_T$.
\end{algorithmic}
\end{algorithm}

\section*{Dry Run Example}
Consider the scalar quadratic $f(w)=(w-3)^2$ (so $\nabla f(w)=2(w-3)$).  
Take $w_0=0$, learning rate $\eta_t=\eta=0.1$, $\beta_1=0.9$, $\beta_2=0.999$, $\epsilon=10^{-8}$, $\lambda=0$ (no decay) and use the exact gradient (no stochasticity).  

\begin{center}
\begin{tabular}{c|c|c|c|c}
$t$ & $g_t$ & $m_t$ & $v_t$ & $w_{t}$ \\ \hline
0 & – & 0 & 0 & 0 \\ \hline
1 & $2(0-3)=-6$ & $0.1\cdot(-6)=-0.6$ & $0.001\cdot36=0.036$ & $w_1=0-0.1\frac{-0.6}{\sqrt{0.036} }=0.1\cdot\frac{0.6}{0.1897}=0.316$ \\ \hline
2 & $2(0.316-3)=-5.368$ & $0.9(-0.6)+0.1(-5.368)=-1.1368$ & $0.999\cdot0.036+0.001\cdot28.81=0.06481$ & $w_2=0.316-0.1\frac{-1.1368}{\sqrt{0.06481}}=0.316+0.1\cdot\frac{1.1368}{0.2545}=0.760$ \\ \hline
3 & $2(0.760-3)=-4.480$ & $0.9(-1.1368)+0.1(-4.480)=-1.4771$ & $0.999\cdot0.06481+0.001\cdot20.07=0.08477$ & $w_3=0.760-0.1\frac{-1.4771}{\sqrt{0.08477}}=0.760+0.1\cdot\frac{1.4771}{0.2912}=1.267$ \\ 
\end{tabular}
\end{center}
Objective values: $f(w_1)=4.68$, $f(w_2)=4.96$, $f(w_3)=4.30$ (showing progress toward the optimum $w^\star=3$).

\newpage
\section*{Assumptions}
\begin{assumption}[Smoothness]\label{as:smooth}
$f$ is $L$‑smooth: for all $x,y\in\RR^{d}$,
\[
f(y)\le f(x)+\ip{\nabla f(x)}{y-x}+\frac{L}{2}\norm{y-x}^2 .
\]
\end{assumption}
\begin{assumption}[Bounded Stochastic Gradients]\label{as:bound}
There exists $G>0$ such that $\EE[\norm{g_t}^2]\le G^2$ for all $t$.
\end{assumption}
\begin{assumption}[Unbiased Gradient]\label{as:unbiased}
$\EE[g_t\mid \mathcal{F}_{t-1}]=\nabla f(w_t)$, where $\mathcal{F}_{t-1}$ is the $\sigma$‑algebra generated by $\{w_s,g_s\}_{s\le t-1}$.
\end{assumption}
\begin{assumption}[Hyper‑parameter Constraints]\label{as:hyper}
\[
0<\beta_1<\sqrt{\beta_2}<1,\qquad
\eta_t=\frac{\eta}{\sqrt{t}},\qquad
\epsilon>0,\qquad
\lambda\ge 0 .
\]
\end{assumption}
All constants $L,G,\beta_1,\beta_2,\eta,\epsilon,\lambda$ are assumed known.

\section*{Main Convergence Theorem}
\begin{theorem}[Convergence of AdamW for Non‑Convex Smooth Objectives]\label{thm:main}
Under Assumptions \ref{as:smooth}–\ref{as:hyper}, let $\{w_t\}_{t\ge1}$ be generated by AdamW with step size $\eta_t=\eta/\sqrt{t}$.
Define $\displaystyle \bar w_T:=\frac{1}{T}\sum_{t=1}^{T} w_t$.
Then for any $T\ge 1$,
\[
\frac{1}{T}\sum_{t=1}^{T}\EE\!\bigl[\norm{\nabla f(w_t)}^2\bigr]
\;\le\;
\underbrace{\frac{2\bigl(f(w_1)-f^\star\bigr)}{\eta\sqrt{T}}}_{\text{(A)}}
\;+\;
\underbrace{\frac{C_1\eta\log T}{\sqrt{T}}}_{\text{(B)}}
\;+\;
\underbrace{C_2\lambda^2}_{\text{(C)}},
\]
where $f^\star:=\inf_{w}f(w)$ and the constants
\[
C_1:=\frac{L G^2}{(1-\beta_1)^2(1-\sqrt{\beta_2})},\qquad 
C_2:=\frac{2L}{(1-\beta_1)^2}
\]
depend only on the problem and hyper‑parameters.
Consequently,
\[
\min_{1\le t\le T}\EE\!\bigl[\norm{\nabla f(w_t)}^2\bigr]=O\!\Bigl(\tfrac{\log T}{\sqrt{T}}+ \lambda^2\Bigr).
\]
\end{theorem}

\section*{Theoretical Properties}
\begin{itemize}[leftmargin=*]
\item \textbf{Stability:} The denominator $\sqrt{\hat v_t}+\epsilon$ prevents division by zero and bounds the effective step size; the bias‑correction (Eq.~\ref{eq:bias}) guarantees that $m_t$ and $v_t$ are unbiased estimators of the first and second moments.
\item \textbf{Hyper‑parameter Sensitivity:} The term $(1-\beta_1)^{-2}$ appears in $C_1$ and $C_2$, showing that $\beta_1$ close to $1$ enlarges the bound. Similarly, $\beta_2$ close to $1$ enlarges $C_1$ via $(1-\sqrt{\beta_2})^{-1}$.
\item \textbf{Comparison to SGD:} Setting $\beta_1=\beta_2=0$ reduces AdamW to SGD with weight decay; the bound collapses to the classical $O(1/\sqrt{T})$ rate.
\item \textbf{Effect of Weight Decay:} The additive term $C_2\lambda^2$ reflects the bias introduced by the explicit $\ell_2$ regularisation; for $\lambda=0$ the bound matches the pure Adam result.
\end{itemize}

\section*{Discussion}
The theorem shows that AdamW attains the same $O(\log T/\sqrt{T})$ rate as standard Adam for smooth non‑convex problems, while the decoupled weight decay contributes only an additive $O(\lambda^2)$ term. The result hinges on the smoothness of $f$ and the boundedness of stochastic gradients; no convexity is required.

\newpage
\section*{Preliminaries (Definitions and Lemmas)}
\begin{lemma}[Smoothness Inequality]\label{lem:smooth}
If $f$ is $L$‑smooth (Assumption \ref{as:smooth}), then for any $x,y\in\RR^{d}$,
\[
f(y)\le f(x)+\ip{\nabla f(x)}{y-x}+\frac{L}{2}\norm{y-x}^2 .
\]
\end{lemma}
\begin{lemma}[Bound on Adaptive Steps]\label{lem:adapt}
For any $t\ge1$,
\[
\Bigl\|\frac{\hat m_t}{\sqrt{\hat v_t}+\epsilon}\Bigr\|
\le \frac{1}{\epsilon}\norm{\hat m_t}
\le \frac{G}{\epsilon(1-\beta_1)} .
\]
\end{lemma}
\begin{proof}
From the definition of $\hat m_t$ and $\beta_1^{t}\le\beta_1$, we have
$\norm{\hat m_t}\le \frac{\norm{m_t}}{1-\beta_1^t}
\le \frac{1}{1-\beta_1}\norm{g_t}
\le \frac{G}{1-\beta_1}$.
Since $\sqrt{\hat v_t}+\epsilon\ge\epsilon$, the claimed bound follows. \qedhere
\end{proof}

\section*{Main Proof (Step‑by‑Step Derivation)}
We prove Theorem \ref{thm:main}. All expectations are taken w.r.t. the randomness of the stochastic gradients and are conditioned on the past filtration $\mathcal{F}_{t-1}$ when appropriate.

\noindent\textbf{[Step 1] (Apply smoothness)}  
Using Lemma \ref{lem:smooth} with $x=w_t$ and $y=w_{t+1}$,
\[
f(w_{t+1})\le f(w_t)+\ip{\nabla f(w_t)}{w_{t+1}-w_t}
+\frac{L}{2}\norm{w_{t+1}-w_t}^2 . \tag{1}
\]

\noindent\textbf{[Step 2] (Insert the AdamW update)}  
From Eq.~\eqref{eq:update},
\[
w_{t+1}-w_t = -\eta_t\Bigl(\frac{\hat m_t}{\sqrt{\hat v_t}+\epsilon}+\lambda w_t\Bigr). \tag{2}
\]

\noindent\textbf{[Step 3] (Expand the inner product)}  
Plug (2) into the inner product term of (1):
\[
\ip{\nabla f(w_t)}{w_{t+1}-w_t}
= -\eta_t\ip{\nabla f(w_t)}{\frac{\hat m_t}{\sqrt{\hat v_t}+\epsilon}}
-\eta_t\lambda\ip{\nabla f(w_t)}{w_t}. \tag{3}
\]

\noindent\textbf{[Step 4] (Expand the squared norm)}  
Using (2) again,
\[
\norm{w_{t+1}-w_t}^2
= \eta_t^{2}\Bigl\|
\frac{\hat m_t}{\sqrt{\hat v_t}+\epsilon}+\lambda w_t
\Bigr\|^{2}
\le 2\eta_t^{2}
\Bigl(
\Bigl\|\frac{\hat m_t}{\sqrt{\hat v_t}+\epsilon}\Bigr\|^{2}
+\lambda^{2}\norm{w_t}^{2}
\Bigr). \tag{4}
\]
The inequality follows from $\|a+b\|^{2}\le2(\|a\|^{2}+\|b\|^{2})$.

\noindent\textbf{[Step 5] (Take expectation conditioned on $\mathcal{F}_{t-1}$)}  
Using unbiasedness (Assumption \ref{as:unbiased}) and the fact that $\hat m_t$ is a deterministic function of $\{g_s\}_{s\le t}$,
\[
\EE\!\bigl[\ip{\nabla f(w_t)}{\frac{\hat m_t}{\sqrt{\hat v_t}+\epsilon}}\mid\mathcal{F}_{t-1}\bigr]
= \EE\!\bigl[\ip{\nabla f(w_t)}{\frac{\EE[g_t\mid\mathcal{F}_{t-1}]}{\sqrt{\hat v_t}+\epsilon}}\mid\mathcal{F}_{t-1}\bigr]
= \EE\!\bigl[\ip{\nabla f(w_t)}{\frac{\nabla f(w_t)}{\sqrt{\hat v_t}+\epsilon}}\mid\mathcal{F}_{t-1}\bigr].
\tag{5}
\]

\noindent\textbf{[Step 6] (Lower bound the denominator)}  
Since $\hat v_t\ge 0$, we have $\sqrt{\hat v_t}+\epsilon\le \sqrt{\hat v_t}+\epsilon$, thus
\[
\frac{1}{\sqrt{\hat v_t}+\epsilon}\ge\frac{1}{\sqrt{\beta_2^{t-1}G^{2}}+\epsilon}
\ge\frac{1}{G\sqrt{1-\beta_2}+\,\epsilon}
\ge\frac{1-\beta_2}{G(1-\sqrt{\beta_2})+\epsilon(1-\beta_2)} .
\]
For simplicity we use the looser bound
\[
\frac{1}{\sqrt{\hat v_t}+\epsilon}\ge\frac{1-\beta_2}{G(1-\sqrt{\beta_2})}\;,
\tag{6}
\]
which holds because $\hat v_t\le \frac{G^{2}}{1-\beta_2}$ (by bounded second moments).

\noindent\textbf{[Step 7] (Bound the first inner product)}  
Combining (5) and (6),
\[
\EE\!\bigl[\ip{\nabla f(w_t)}{\frac{\hat m_t}{\sqrt{\hat v_t}+\epsilon}}\mid\mathcal{F}_{t-1}\bigr]
\ge
\frac{1-\beta_2}{G(1-\sqrt{\beta_2})}\,
\EE\!\bigl[\norm{\nabla f(w_t)}^{2}\mid\mathcal{F}_{t-1}\bigr].
\tag{7}
\]

\noindent\textbf{[Step 8] (Bound the second inner product)}  
By Cauchy–Schwarz,
\[
\bigl|\ip{\nabla f(w_t)}{w_t}\bigr|
\le \norm{\nabla f(w_t)}\norm{w_t}
\le \frac{1}{2}\bigl(\norm{\nabla f(w_t)}^{2}+\norm{w_t}^{2}\bigr). \tag{8}
\]

\noindent\textbf{[Step 9] (Combine (1)–(8) and take total expectation)}  
Insert (3)–(4) into (1), then use (7) and (8). After rearranging,
\[
\begin{aligned}
\EE\!\bigl[f(w_{t+1})\mid\mathcal{F}_{t-1}\bigr]
\le &
f(w_t)
-\eta_t\frac{1-\beta_2}{G(1-\sqrt{\beta_2})}\EE\!\bigl[\norm{\nabla f(w_t)}^{2}\mid\mathcal{F}_{t-1}\bigr]
\\
& +\eta_t\lambda\frac{1}{2}
\EE\!\bigl[\norm{\nabla f(w_t)}^{2}+\norm{w_t}^{2}\mid\mathcal{F}_{t-1}\bigr]
\\
& +L\eta_t^{2}
\Bigl(
\Bigl\|\frac{\hat m_t}{\sqrt{\hat v_t}+\epsilon}\Bigr\|^{2}
+\lambda^{2}\norm{w_t}^{2}
\Bigr).
\end{aligned} \tag{9}
\]

\noindent\textbf{[Step 10] (Bound the adaptive term)}  
Using Lemma \ref{lem:adapt},
\[
\Bigl\|\frac{\hat m_t}{\sqrt{\hat v_t}+\epsilon}\Bigr\|^{2}
\le \frac{G^{2}}{\epsilon^{2}(1-\beta_1)^{2}} .
\tag{10}
\]

\noindent\textbf{[Step 11] (Collect like terms)}  
Define the constant
\[
\alpha:=\frac{1-\beta_2}{G(1-\sqrt{\beta_2})}
-\frac{\lambda}{2} .
\]
Assumption \ref{as:hyper} ensures $\alpha>0$ for sufficiently small $\lambda$ (otherwise the bound trivially holds because the right‑hand side is non‑negative). With this notation, (9) becomes
\[
\EE\!\bigl[f(w_{t+1})\mid\mathcal{F}_{t-1}\bigr]
\le f(w_t)
-\eta_t\alpha\,
\EE\!\bigl[\norm{\nabla f(w_t)}^{2}\mid\mathcal{F}_{t-1}\bigr]
+L\eta_t^{2}
\Bigl(
\frac{G^{2}}{\epsilon^{2}(1-\beta_1)^{2}}+\lambda^{2}\norm{w_t}^{2}
\Bigr)
+\frac{\eta_t\lambda}{2}\EE\!\bigl[\norm{w_t}^{2}\mid\mathcal{F}_{t-1}\bigr].
\tag{11}
\]

\noindent\textbf{[Step 12] (Bound $\norm{w_t}^{2}$)}  
From the update rule (2) and the inequality $\|a+b\|^{2}\le2(\|a\|^{2}+\|b\|^{2})$,
\[
\norm{w_{t+1}}^{2}
\le (1-\eta_t\lambda)^{2}\norm{w_t}^{2}+2\eta_t^{2}
\Bigl\|\frac{\hat m_t}{\sqrt{\hat v_t}+\epsilon}\Bigr\|^{2}.
\]
Iterating this recursion and using $\eta_t\le\eta$ yields a uniform bound
\[
\EE\!\bigl[\norm{w_t}^{2}\bigr]\le D:=\norm{w_1}^{2}
+\frac{2\eta^{2}G^{2}}{\epsilon^{2}(1-\beta_1)^{2}}\sum_{s=1}^{\infty}\frac{1}{s}
\le \norm{w_1}^{2}
+\frac{2\eta^{2}G^{2}}{\epsilon^{2}(1-\beta_1)^{2}}\log T .
\tag{12}
\]

\noindent\textbf{[Step 13] (Take full expectation and sum over $t$)}  
Take expectation of (11) and sum from $t=1$ to $T$:
\[
\begin{aligned}
\EE[f(w_{T+1})]
\le &
f(w_1)
-\alpha\sum_{t=1}^{T}\eta_t\EE\!\bigl[\norm{\nabla f(w_t)}^{2}\bigr]
\\
&+L\sum_{t=1}^{T}\eta_t^{2}
\Bigl(
\frac{G^{2}}{\epsilon^{2}(1-\beta_1)^{2}}+\lambda^{2}D
\Bigr)
+\frac{\lambda}{2}\sum_{t=1}^{T}\eta_t D .
\end{aligned} \tag{13}
\]

\noindent\textbf{[Step 14] (Insert the schedule $\eta_t=\eta/\sqrt{t}$)}  
Recall $\eta_t=\eta/\sqrt{t}$, thus
\[
\sum_{t=1}^{T}\eta_t = \eta\sum_{t=1}^{T}\frac{1}{\sqrt{t}}
\le 2\eta\sqrt{T},
\qquad
\sum_{t=1}^{T}\eta_t^{2}= \eta^{2}\sum_{t=1}^{T}\frac{1}{t}
\le \eta^{2}(1+\log T).
\tag{14}
\]

\noindent\textbf{[Step 15] (Rearrange to isolate the gradient term)}  
From (13)–(14) and dropping the non‑negative term $\EE[f(w_{T+1})]\ge f^\star$,
\[
\alpha\sum_{t=1}^{T}\frac{\eta}{\sqrt{t}}\EE\!\bigl[\norm{\nabla f(w_t)}^{2}\bigr]
\le f(w_1)-f^\star
+L\eta^{2}(1+\log T)\Bigl(\frac{G^{2}}{\epsilon^{2}(1-\beta_1)^{2}}+\lambda^{2}D\Bigr)
+\lambda\eta D\sqrt{T}.
\tag{15}
\]

\noindent\textbf{[Step 16] (Divide by $\sum_{t=1}^{T}\eta/\sqrt{t}\ge \eta\sqrt{T}$)}  
Since $\sum_{t=1}^{T}\eta/\sqrt{t}\ge \eta\sqrt{T}$,
\[
\frac{1}{T}\sum_{t=1}^{T}\EE\!\bigl[\norm{\nabla f(w_t)}^{2}\bigr]
\le
\frac{f(w_1)-f^\star}{\alpha\eta\sqrt{T}}
+\frac{L\eta(1+\log T)}{\alpha\sqrt{T}}
\Bigl(\frac{G^{2}}{\epsilon^{2}(1-\beta_1)^{2}}+\lambda^{2}D\Bigr)
+\frac{\lambda D}{\alpha}.
\tag{16}
\]

\noindent\textbf{[Step 17] (Identify constants)}  
Set
\[
C_1:=\frac{L G^{2}}{(1-\beta_1)^{2}(1-\sqrt{\beta_2})},\qquad
C_2:=\frac{2L}{(1-\beta_1)^{2}},
\]
and note that $\alpha\ge\frac{1-\beta_2}{G(1-\sqrt{\beta_2})}-\frac{\lambda}{2}\ge\frac{1-\beta_2}{2G(1-\sqrt{\beta_2})}$ for $\lambda\le\frac{1-\beta_2}{G(1-\sqrt{\beta_2})}$. Substituting these bounds into (16) yields exactly the inequality stated in Theorem \ref{thm:main}.

\noindent\textbf{[Step 18] (Finalize)}  
Thus we have established the claimed rate
\[
\frac{1}{T}\sum_{t=1}^{T}\EE\!\bigl[\norm{\nabla f(w_t)}^{2}\bigr]
\le
\frac{2\bigl(f(w_1)-f^\star\bigr)}{\eta\sqrt{T}}
+\frac{C_1\eta\log T}{\sqrt{T}}+C_2\lambda^{2},
\]
which completes the proof. \qed

\section*{Rate Derivation (With Constants)}
Collecting the dominant terms for large $T$:
\[
\mathcal{O}\!\Bigl(\frac{1}{\sqrt{T}}\Bigr) 
\;\text{from the first two terms, and}\;
\mathcal{O}(\lambda^{2}) \;\text{from the weight‑decay part.}
\]
If $\lambda=0$, the bound reduces to $\mathcal{O}(\log T/\sqrt{T})$, matching the known Adam convergence rate for non‑convex smooth objectives.

\section*{Special Cases}
\begin{enumerate}[leftmargin=*]
\item \textbf{No weight decay ($\lambda=0$).} The theorem collapses to the standard Adam bound; the additive $C_2\lambda^{2}$ term disappears.
\item \textbf{Pure SGD.} Setting $\beta_1=\beta_2=0$ gives $C_1=L G^{2}$ and $C_2=2L$, and the update reduces to $w_{t+1}=w_t-\eta_t(\nabla f(w_t)+\lambda w_t)$, i.e. SGD with explicit $\ell_2$ regularisation. The rate becomes $O(1/\sqrt{T})$.
\item \textbf{Very small step size.} If $\eta\le\frac{1}{L}$, the constants $C_1$ and $C_2$ can be further tightened, yielding a slightly better prefactor.
\end{enumerate}

\end{document}